% !TEX program = pdflatex
\documentclass[11pt,a4paper]{article}

% =========================================================================
% PLAN DEL EQUIPO
% Proyecto: CasaScore (AlquilaPuno)
% Gestión de Proyectos TICS (EST315) -- Semestre 2026-II -- Team Sobrevivientes
% =========================================================================

% --- PAQUETES DE CONFIGURACIÓN ---
\usepackage[utf8]{inputenc}
\usepackage[T1]{fontenc}
\usepackage[spanish,es-tabla]{babel}
% Compatibilidad con babel-spanish reciente (flechas de TikZ y signo de porcentaje)
\AtBeginDocument{\shorthandoff{<>}\spanishplainpercent}
\usepackage{geometry}
\geometry{
    top=3.6cm,
    bottom=2.3cm,
    left=2.2cm,
    right=2.2cm,
    % PARTE MODIFICABLE: headheight debe ser mayor que la altura del logo más alto (2.56cm = 73pt).
    headheight=80pt,
    headsep=15pt,
    footskip=20pt
}

\usepackage{graphicx}
\usepackage{fancyhdr}
\usepackage[table]{xcolor}
\usepackage{titlesec}
\usepackage{booktabs}
\usepackage{tabularx}
\usepackage{array}
\usepackage{amsmath}
\usepackage{enumitem}
\usepackage{microtype}
\usepackage{amssymb}
% --- RECURSOS VISUALES: GRÁFICOS Y CÓDIGO DE COLORES ---
\usepackage{tikz}
\usetikzlibrary{positioning,calc}
\definecolor{GrisFila}{RGB}{246, 248, 251}
\definecolor{CuadCerca}{RGB}{150, 30, 30}
\definecolor{CuadSatisf}{RGB}{180, 130, 20}
\definecolor{CuadInfo}{RGB}{10, 45, 90}
\definecolor{CuadMonit}{RGB}{110, 115, 125}
\newcommand{\etiqueta}[2]{{\setlength{\fboxsep}{2pt}\colorbox{#1}{\textcolor{white}{\scriptsize\bfseries #2}}}}
\newcommand{\cerca}{\etiqueta{CuadCerca}{Gestionar de cerca}}
\newcommand{\satisf}{\etiqueta{CuadSatisf}{Mantener satisfecho}}
\newcommand{\informado}{\etiqueta{CuadInfo}{Mantener informado}}
\newcommand{\monit}{\etiqueta{CuadMonit}{Monitorear}}
\usepackage{needspace}
\usepackage{longtable}
\definecolor{Rama1}{RGB}{10, 45, 90}
\definecolor{Rama2}{RGB}{40, 95, 150}
\definecolor{Rama3}{RGB}{180, 130, 20}
\definecolor{Rama4}{RGB}{30, 110, 60}
\definecolor{Rama5}{RGB}{110, 60, 140}
\definecolor{RaciR}{RGB}{150, 30, 30}
\definecolor{RaciA}{RGB}{10, 45, 90}
\definecolor{RaciC}{RGB}{180, 130, 20}
\definecolor{RaciI}{RGB}{140, 145, 155}
\definecolor{NivAlto}{RGB}{30, 110, 60}
\definecolor{NivMedio}{RGB}{180, 130, 20}
\definecolor{NivBajo}{RGB}{150, 30, 30}

\usepackage{hyperref}

% Logos: funciona a 2 o 3 niveles de profundidad dentro de docs/latex/
\graphicspath{{../../assets/img/}{../../../assets/img/}}

% --- PALETA INSTITUCIONAL FIEI / UNA PUNO ---
\definecolor{AzulUNA}{RGB}{10, 45, 90}
\definecolor{DoradoFIEI}{RGB}{180, 130, 20}
\definecolor{GrisOscuro}{RGB}{50, 50, 50}
\definecolor{AzulClaro}{RGB}{228, 236, 246}

% --- CONFIGURACIÓN DE LISTAS ---
\setlist{topsep=2pt, itemsep=2pt, parsep=1pt}

% --- COLUMNAS DE TABLA ---
\newcolumntype{L}{>{\raggedright\arraybackslash}X}
\newcolumntype{P}[1]{>{\raggedright\arraybackslash}p{#1}}
\newcolumntype{C}[1]{>{\centering\arraybackslash}p{#1}}
\newcolumntype{Q}[1]{>{\raggedleft\arraybackslash}p{#1}}
\newcolumntype{R}{>{\raggedleft\arraybackslash}X}
\renewcommand{\arraystretch}{1.25}

% --- FORMATO DE ENLACES ---
\hypersetup{
    colorlinks=true,
    linkcolor=AzulUNA,
    urlcolor=AzulUNA,
    citecolor=AzulUNA,
    pdftitle={PLAN DEL EQUIPO -- CasaScore (AlquilaPuno)},
    pdfauthor={Team Sobrevivientes}
}

% --- FORMATO DE TÍTULOS ---
\titleformat{\section}{\color{AzulUNA}\normalfont\large\bfseries}{\thesection.}{0.7em}{}[\color{AzulUNA}\titlerule]
\titleformat{\subsection}{\color{AzulUNA}\normalfont\normalsize\bfseries}{\thesubsection.}{0.5em}{}
\titlespacing*{\section}{0pt}{9pt}{4pt}
\titlespacing*{\subsection}{0pt}{6pt}{2pt}

% --- ENCABEZADO Y PIE DE PÁGINA (fancyhdr) ---
\pagestyle{fancy}
\fancyhf{}

% ===== TAMAÑO DE LOS LOGOS (PARTE MODIFICABLE) =====
% logo_fiei.png  : height=1.45cm (lado izquierdo)
% logo_equipo.png: height=2.56cm (lado derecho; 25 % más grande porque la imagen tiene más borde blanco)
% Para cambiar el tamaño, modifique solo el valor height=... de cada \includegraphics.
% Si el logo más alto supera 2.6cm, aumente también headheight y top en la configuración de página.
\fancyhead[L]{%
    \includegraphics[height=1.45cm,keepaspectratio]{logo_fiei.png}%
}
\fancyhead[C]{%
    \scriptsize \textbf{\textcolor{AzulUNA}{UNIVERSIDAD NACIONAL DEL ALTIPLANO DE PUNO}}\\[2pt]
    \tiny \textcolor{GrisOscuro}{FACULTAD DE INGENIERÍA ESTADÍSTICA E INFORMÁTICA}\\[2pt]
    \scriptsize \textbf{\textcolor{AzulUNA}{Gestión de Proyectos TICS (EST315) \textbar\ Team Sobrevivientes}}\\[2pt]
    \footnotesize \textbf{\textcolor{DoradoFIEI}{PLAN DEL EQUIPO}}%
}
\fancyhead[R]{%
    \includegraphics[height=2.56cm,keepaspectratio]{logo_equipo.png}%
}

\renewcommand{\headrulewidth}{0pt}
\renewcommand{\headrule}{}

\fancyfoot[L]{\scriptsize \textcolor{GrisOscuro}{CasaScore (AlquilaPuno)}}
\fancyfoot[C]{\small \thepage}
\fancyfoot[R]{\scriptsize \textcolor{GrisOscuro}{Puno, Perú -- Semestre 2026-II}}
\renewcommand{\footrulewidth}{0.4pt}

\begin{document}

\begin{center}
    {\Large \textbf{\textcolor{AzulUNA}{PLAN DEL EQUIPO}}}\\[3pt]
    {\normalsize \textbf{CasaScore (AlquilaPuno)}}\\[4pt]
    \footnotesize \textbf{Fecha:} viernes 18/09/2026 (semana 4) \quad \textbar\quad \textbf{Versión:} 1.0 \quad \textbar\quad \textbf{Fase:} Planificación
\end{center}

\vspace{0.1cm}

\section{Datos Generales}
\noindent\begin{tabularx}{\textwidth}{@{}P{4.2cm} L@{}}
\toprule
\textbf{Proyecto:} & CasaScore (AlquilaPuno) \\
\textbf{Responsable:} & Ruth Karina Apaza Solis -- Analista de Planificación \\
\textbf{Proceso PM4DEV:} & Gestión del equipo del proyecto \\
\textbf{Insumos:} & Acta de conformación del equipo, Acta de Constitución y EDT \\
\textbf{Periodo:} & Semanas 2 a 17 del semestre 2026-II (16 semanas) \\
\bottomrule
\end{tabularx}

\section{Organigrama}
El equipo tiene una estructura plana: el Project Manager coordina a tres roles especializados y el docente del curso acompaña como asesor, fuera de la línea de mando.

\begin{center}
\begin{tikzpicture}[font=\footnotesize,
  caja/.style={draw=#1, fill=#1!8, line width=0.8pt, rounded corners=3pt, text width=3.45cm, align=center, inner sep=4pt},
  rol/.style={font=\scriptsize\bfseries, text=#1}]
\node[caja=AzulUNA, text width=4.2cm] (pm) at (0,0) {{\scriptsize\bfseries\textcolor{AzulUNA}{PROJECT MANAGER}}\\ Clyde Neil Paricahua Pari\\ {\tiny Dirige, decide y responde por el proyecto}};
\node[caja=GrisOscuro, dashed, text width=2.9cm] (as) at (5.0,0) {{\scriptsize\bfseries\textcolor{GrisOscuro}{ASESOR}}\\ Ing. Renzo Apaza Cutipa\\ {\tiny Orienta; visto bueno en el acta de constitución y de cierre}};
\node[caja=Rama2] (ru) at (-4.6,-2.4) {{\scriptsize\bfseries\textcolor{Rama2}{ANALISTA DE PLANIFICACIÓN}}\\ Ruth Karina Apaza Solis\\ {\tiny Alcance, tiempo, costo y control}};
\node[caja=Rama3] (el) at (0,-2.4) {{\scriptsize\bfseries\textcolor{Rama3}{ESPECIALISTA TÉCNICO}}\\ Eliseo Tarqui Ajahuana\\ {\tiny Arquitectura, desarrollo y pruebas técnicas}};
\node[caja=Rama4] (vi) at (4.6,-2.4) {{\scriptsize\bfseries\textcolor{Rama4}{GESTORA DE COMUNICACIÓN}}\\ Viviana Mary Quisbert Quispe\\ {\tiny Riesgos, calidad, actas y documentación}};
\draw[AzulUNA, thick] (pm.south) -- ++(0,-0.45) coordinate (m);
\draw[AzulUNA, thick] (m) -| (ru.north); \draw[AzulUNA, thick] (m) -- (el.north); \draw[AzulUNA, thick] (m) -| (vi.north);
\draw[GrisOscuro, dashed, thick, <->] (pm.east) -- node[above, font=\tiny]{asesoría} (as.west);
\node[draw=AzulUNA!50, fill=AzulClaro!60, rounded corners=3pt, minimum width=12.6cm, text width=12.4cm, align=center, minimum height=0.55cm, font=\scriptsize] at (0,-4.55) {\textbf{Team Sobrevivientes}, sponsor y dueño del proyecto: los cuatro integrantes aprueban los planes y los cambios mayores por unanimidad.};
\end{tikzpicture}
\end{center}

\Needspace{28\baselineskip}
\section{Roles y Responsabilidades}
{\small
\setlength{\tabcolsep}{3.5pt}
\rowcolors{2}{white}{GrisFila}
\noindent\begin{tabularx}{\textwidth}{@{}P{2.6cm} P{3.2cm} L P{1.9cm} P{3.1cm}@{}}
\toprule
\rowcolor{AzulClaro}
\textbf{Rol} & \textbf{Propósito} & \textbf{Responsabilidades} & \textbf{Reporta a} & \textbf{Cómo se evalúa} \\
\midrule
\textbf{Project Manager}\newline {\footnotesize Clyde Neil Paricahua Pari} & Dirigir el proyecto para cumplir alcance, tiempo, costo y satisfacción del cliente. & Conducir las reuniones; aprobar entregables como responsable final; gestionar interesados externos y la relación con el asesor; decidir sobre cambios menores. & Equipo (sponsor) y asesor & Entregables aprobados a tiempo y decisiones registradas en actas. \\
\textbf{Analista de Planificación}\newline {\footnotesize Ruth Karina Apaza Solis} & Mantener las líneas base de alcance, tiempo y costo y medir el avance. & EDT, cronograma y presupuesto; registro de horas y gastos; valor ganado y tablero de control semanal. & Project Manager & Informe de avance cada martes y desviaciones de costo y plazo menores a 10\,\%. \\
\textbf{Especialista Técnico}\newline {\footnotesize Eliseo Tarqui Ajahuana} & Construir un producto que funcione y cumpla los requerimientos. & Arquitectura, base de datos y algoritmos; desarrollo de las funciones; revisión técnica de los Pull Request; manuales. & Project Manager & Funciones terminadas por sprint y 0 errores críticos en la entrega. \\
\textbf{Gestora de Comunicación}\newline {\footnotesize Viviana Mary Quisbert Quispe} & Asegurar la calidad, controlar los riesgos y mantener informados a los interesados. & Políticas y actas; registro de riesgos; plan de calidad y pruebas; comunicaciones y expediente en \LaTeX; lecciones aprendidas. & Project Manager & Actas al día, riesgos revisados cada semana y casos de prueba ejecutados. \\
\bottomrule
\end{tabularx}
}

\section{Tabla de Habilidades}
La tabla compara las habilidades que exige cada rama de la EDT con el nivel de cada integrante. Las celdas \etiqueta{NivBajo}{Básica} en una rama muestran dónde el integrante necesita apoyo o capacitación.

\vspace{0.1cm}
{\small
\setlength{\tabcolsep}{3.5pt}
\rowcolors{2}{white}{GrisFila}
\noindent\begin{tabularx}{\textwidth}{@{}P{3.4cm} L C{1.35cm} C{1.35cm} C{1.35cm} C{1.35cm}@{}}
\toprule
\rowcolor{AzulClaro}
\textbf{Rama de la EDT} & \textbf{Habilidad requerida} & \textbf{Clyde} & \textbf{Ruth} & \textbf{Eliseo} & \textbf{Viviana} \\
\midrule
\textcolor{Rama1}{\textbf{1}} Gestión del proyecto & Gestión, planificación y costos & \etiqueta{NivAlto}{Alta} & \etiqueta{NivAlto}{Alta} & \etiqueta{NivBajo}{Básica} & \etiqueta{NivMedio}{Media} \\
\textcolor{Rama2}{\textbf{2}} Análisis y diseño & Requerimientos, modelo de datos y bocetos & \etiqueta{NivMedio}{Media} & \etiqueta{NivMedio}{Media} & \etiqueta{NivAlto}{Alta} & \etiqueta{NivBajo}{Básica} \\
\textcolor{Rama3}{\textbf{3}} Construcción del software & Expo, React Native, TypeScript, SQLite y Supabase & \etiqueta{NivBajo}{Básica} & \etiqueta{NivBajo}{Básica} & \etiqueta{NivAlto}{Alta} & \etiqueta{NivBajo}{Básica} \\
\textcolor{Rama4}{\textbf{4}} Pruebas y validación & Casos de prueba, encuestas y trato con clientes & \etiqueta{NivAlto}{Alta} & \etiqueta{NivMedio}{Media} & \etiqueta{NivMedio}{Media} & \etiqueta{NivAlto}{Alta} \\
\textcolor{Rama5}{\textbf{5}} Cierre & Redacción técnica, \LaTeX\ y exposición & \etiqueta{NivAlto}{Alta} & \etiqueta{NivMedio}{Media} & \etiqueta{NivBajo}{Básica} & \etiqueta{NivAlto}{Alta} \\
\bottomrule
\end{tabularx}
}

\vspace{0.2cm}
\Needspace{8\baselineskip}
\noindent\textbf{Brechas y acciones}

\vspace{0.1cm}
{\small
\rowcolors{2}{white}{GrisFila}
\noindent\begin{tabularx}{\textwidth}{@{}P{4.6cm} L P{2.9cm}@{}}
\toprule
\rowcolor{AzulClaro}
\textbf{Brecha} & \textbf{Acción} & \textbf{Responsable} \\
\midrule
Construcción concentrada en el Especialista Técnico (lidera todos los paquetes de la rama 3). & Los demás integrantes apoyan en pruebas, datos y pantallas sencillas; revisión de código en cada Pull Request. & Eliseo Tarqui Ajahuana \\
Uso de Supabase y seguridad por fila, nuevo para el equipo. & Práctica guiada antes de la semana 9 con la documentación oficial. & Eliseo Tarqui Ajahuana \\
\bottomrule
\end{tabularx}
}

\Needspace{14\baselineskip}
\section{Matriz de Asignación}
Todos los integrantes dedican la misma carga: \textbf{8 horas por semana} durante 16 semanas, es decir 128 horas cada uno y 512 horas del equipo, igual al total de la EDT. Sumadas a las 8 horas semanales de Horus Sentinel, cada integrante trabaja 16 horas por semana en la cartera del equipo.

\vspace{0.1cm}
{\small
\setlength{\tabcolsep}{4pt}
\noindent\begin{tabularx}{\textwidth}{@{}L C{1.0cm} C{1.0cm} C{1.0cm} C{1.0cm} C{1.0cm} C{1.1cm} C{1.1cm}@{}}
\toprule
\rowcolor{AzulClaro}
\textbf{Integrante} & \textcolor{Rama1}{\textbf{Rama 1}} & \textcolor{Rama2}{\textbf{Rama 2}} & \textcolor{Rama3}{\textbf{Rama 3}} & \textcolor{Rama4}{\textbf{Rama 4}} & \textcolor{Rama5}{\textbf{Rama 5}} & \textbf{Total h} & \textbf{h/sem.} \\
\midrule
Clyde Neil Paricahua Pari & 56 & 8 & 24 & 24 & 16 & \textbf{128} & \textbf{8} \\
Ruth Karina Apaza Solis & 56 & 16 & 24 & 24 & 8 & \textbf{128} & \textbf{8} \\
Eliseo Tarqui Ajahuana & 16 & 32 & 64 & 8 & 8 & \textbf{128} & \textbf{8} \\
Viviana Mary Quisbert Quispe & 40 & 8 & 48 & 16 & 16 & \textbf{128} & \textbf{8} \\
\midrule
\textbf{Total} & \textbf{168} & \textbf{64} & \textbf{160} & \textbf{72} & \textbf{48} & \textbf{512} & \textbf{32} \\
\bottomrule
\end{tabularx}
}

\vspace{0.15cm}
\noindent\textbf{Distribución de las horas de cada integrante por rama}
\begin{center}
\begin{tikzpicture}[x=0.085cm, y=0.62cm, font=\scriptsize]
\node[left, text width=4.2cm, align=right] at (-1,0) {Clyde Neil Paricahua Pari};
\fill[Rama1] (0,-0.30) rectangle (56,0.30);
\node[text=white, font=\tiny\bfseries] at (28.0,0) {56};
\fill[Rama2] (56,-0.30) rectangle (64,0.30);
\fill[Rama3] (64,-0.30) rectangle (88,0.30);
\node[text=white, font=\tiny\bfseries] at (76.0,0) {24};
\fill[Rama4] (88,-0.30) rectangle (112,0.30);
\node[text=white, font=\tiny\bfseries] at (100.0,0) {24};
\fill[Rama5] (112,-0.30) rectangle (128,0.30);
\node[text=white, font=\tiny\bfseries] at (120.0,0) {16};
\node[left, text width=4.2cm, align=right] at (-1,-1) {Ruth Karina Apaza Solis};
\fill[Rama1] (0,-1.30) rectangle (56,-0.70);
\node[text=white, font=\tiny\bfseries] at (28.0,-1) {56};
\fill[Rama2] (56,-1.30) rectangle (72,-0.70);
\node[text=white, font=\tiny\bfseries] at (64.0,-1) {16};
\fill[Rama3] (72,-1.30) rectangle (96,-0.70);
\node[text=white, font=\tiny\bfseries] at (84.0,-1) {24};
\fill[Rama4] (96,-1.30) rectangle (120,-0.70);
\node[text=white, font=\tiny\bfseries] at (108.0,-1) {24};
\fill[Rama5] (120,-1.30) rectangle (128,-0.70);
\node[left, text width=4.2cm, align=right] at (-1,-2) {Eliseo Tarqui Ajahuana};
\fill[Rama1] (0,-2.30) rectangle (16,-1.70);
\node[text=white, font=\tiny\bfseries] at (8.0,-2) {16};
\fill[Rama2] (16,-2.30) rectangle (48,-1.70);
\node[text=white, font=\tiny\bfseries] at (32.0,-2) {32};
\fill[Rama3] (48,-2.30) rectangle (112,-1.70);
\node[text=white, font=\tiny\bfseries] at (80.0,-2) {64};
\fill[Rama4] (112,-2.30) rectangle (120,-1.70);
\fill[Rama5] (120,-2.30) rectangle (128,-1.70);
\node[left, text width=4.2cm, align=right] at (-1,-3) {Viviana Mary Quisbert Quispe};
\fill[Rama1] (0,-3.30) rectangle (40,-2.70);
\node[text=white, font=\tiny\bfseries] at (20.0,-3) {40};
\fill[Rama2] (40,-3.30) rectangle (48,-2.70);
\fill[Rama3] (48,-3.30) rectangle (96,-2.70);
\node[text=white, font=\tiny\bfseries] at (72.0,-3) {48};
\fill[Rama4] (96,-3.30) rectangle (112,-2.70);
\node[text=white, font=\tiny\bfseries] at (104.0,-3) {16};
\fill[Rama5] (112,-3.30) rectangle (128,-2.70);
\node[text=white, font=\tiny\bfseries] at (120.0,-3) {16};
\draw[gray!50] (0,-3.45) -- (0,-3.55) node[below, font=\tiny] {0};
\draw[gray!50] (16,-3.45) -- (16,-3.55) node[below, font=\tiny] {16};
\draw[gray!50] (32,-3.45) -- (32,-3.55) node[below, font=\tiny] {32};
\draw[gray!50] (48,-3.45) -- (48,-3.55) node[below, font=\tiny] {48};
\draw[gray!50] (64,-3.45) -- (64,-3.55) node[below, font=\tiny] {64};
\draw[gray!50] (80,-3.45) -- (80,-3.55) node[below, font=\tiny] {80};
\draw[gray!50] (96,-3.45) -- (96,-3.55) node[below, font=\tiny] {96};
\draw[gray!50] (112,-3.45) -- (112,-3.55) node[below, font=\tiny] {112};
\draw[gray!50] (128,-3.45) -- (128,-3.55) node[below, font=\tiny] {128};
\draw[gray!50] (0,-3.5) -- (128,-3.5);
\node[font=\tiny] at (64,-4.25) {horas en el semestre (8 h por semana $\times$ 16 semanas)};
\end{tikzpicture}
\\[2pt]
{\scriptsize \textcolor{Rama1}{\rule{7pt}{7pt}}~Gestión del proyecto \quad \textcolor{Rama2}{\rule{7pt}{7pt}}~Análisis y diseño \quad \textcolor{Rama3}{\rule{7pt}{7pt}}~Construcción del software \quad \textcolor{Rama4}{\rule{7pt}{7pt}}~Pruebas y validación con clientes \quad \textcolor{Rama5}{\rule{7pt}{7pt}}~Cierre}
\end{center}

\Needspace{10\baselineskip}
\section{Matriz RACI}
Cada paquete de trabajo tiene un solo responsable de ejecutarlo y un solo responsable de aprobarlo.

\vspace{0.1cm}
{\small
\noindent\begin{tabularx}{\textwidth}{@{}C{1.2cm} L C{1.2cm} L@{}}
\etiqueta{RaciR}{R} & Responsable: ejecuta y dirige el paquete. & \etiqueta{RaciA}{A} & Aprueba: acepta el entregable y responde por él. \\
\etiqueta{RaciC}{C} & Consultado: aporta antes y durante el trabajo. & \etiqueta{RaciI}{I} & Informado: conoce el resultado. \\
\end{tabularx}
}

\vspace{0.1cm}
{\small
\setlength{\tabcolsep}{3pt}
\begin{longtable}{@{}P{0.75cm} >{\raggedright\arraybackslash}p{5.6cm} C{1.6cm} C{1.6cm} C{1.6cm} C{1.6cm} C{1.6cm}@{}}
\toprule
\rowcolor{AzulClaro}
\textbf{Cód.} & \textbf{Paquete de trabajo} & \textbf{Clyde} & \textbf{Ruth} & \textbf{Eliseo} & \textbf{Viviana} & \textbf{Asesor} \\
\midrule
\endfirsthead
\toprule
\rowcolor{AzulClaro}
\textbf{Cód.} & \textbf{Paquete de trabajo} & \textbf{Clyde} & \textbf{Ruth} & \textbf{Eliseo} & \textbf{Viviana} & \textbf{Asesor} \\
\midrule
\endhead
\bottomrule
\endfoot
\rowcolor{Rama1!15}
\multicolumn{7}{@{}l}{\textbf{\textcolor{Rama1}{1 Gestión del proyecto}}} \\
1.1 & Gobernanza del equipo & \etiqueta{RaciA}{A} & \etiqueta{RaciC}{C} & \etiqueta{RaciI}{I} & \etiqueta{RaciR}{R} & \textcolor{GrisOscuro}{--} \\
1.2 & Viabilidad y Acta de Constitución & \etiqueta{RaciR}{A/R} & \etiqueta{RaciC}{C} & \etiqueta{RaciI}{I} & \etiqueta{RaciC}{C} & \etiqueta{RaciC}{C} \\
1.3 & Planificación base & \etiqueta{RaciA}{A} & \etiqueta{RaciR}{R} & \etiqueta{RaciI}{I} & \etiqueta{RaciC}{C} & \textcolor{GrisOscuro}{--} \\
1.4 & Planes de soporte & \etiqueta{RaciA}{A} & \etiqueta{RaciC}{C} & \etiqueta{RaciI}{I} & \etiqueta{RaciR}{R} & \textcolor{GrisOscuro}{--} \\
1.5 & Seguimiento y control & \etiqueta{RaciA}{A} & \etiqueta{RaciR}{R} & \etiqueta{RaciI}{I} & \etiqueta{RaciC}{C} & \textcolor{GrisOscuro}{--} \\
1.6 & Entregas de la Unidad II & \etiqueta{RaciR}{A/R} & \etiqueta{RaciC}{C} & \etiqueta{RaciI}{I} & \etiqueta{RaciC}{C} & \etiqueta{RaciI}{I} \\
\rowcolor{Rama2!15}
\multicolumn{7}{@{}l}{\textbf{\textcolor{Rama2}{2 Análisis y diseño}}} \\
2.1 & Ficha técnica y requerimientos & \etiqueta{RaciA}{A} & \etiqueta{RaciC}{C} & \etiqueta{RaciR}{R} & \etiqueta{RaciI}{I} & \textcolor{GrisOscuro}{--} \\
2.2 & Modelo de datos y arquitectura & \etiqueta{RaciA}{A} & \etiqueta{RaciC}{C} & \etiqueta{RaciR}{R} & \etiqueta{RaciI}{I} & \textcolor{GrisOscuro}{--} \\
2.3 & Diseño de pantallas & \etiqueta{RaciA}{A} & \etiqueta{RaciC}{C} & \etiqueta{RaciR}{R} & \etiqueta{RaciI}{I} & \textcolor{GrisOscuro}{--} \\
\rowcolor{Rama3!15}
\multicolumn{7}{@{}l}{\textbf{\textcolor{Rama3}{3 Construcción del software}}} \\
3.1 & Motor CASA SCORE & \etiqueta{RaciA}{A} & \etiqueta{RaciI}{I} & \etiqueta{RaciR}{R} & \etiqueta{RaciC}{C} & \textcolor{GrisOscuro}{--} \\
3.2 & Publicación de cuartos & \etiqueta{RaciA}{A} & \etiqueta{RaciI}{I} & \etiqueta{RaciR}{R} & \etiqueta{RaciC}{C} & \textcolor{GrisOscuro}{--} \\
3.3 & Ranking y detalle (MVP Unidad I) & \etiqueta{RaciA}{A} & \etiqueta{RaciI}{I} & \etiqueta{RaciR}{R} & \etiqueta{RaciC}{C} & \textcolor{GrisOscuro}{--} \\
3.4 & Nube, registro e inicio de sesión & \etiqueta{RaciA}{A} & \etiqueta{RaciI}{I} & \etiqueta{RaciR}{R} & \etiqueta{RaciC}{C} & \textcolor{GrisOscuro}{--} \\
3.5 & Punto de referencia y nuevos criterios & \etiqueta{RaciA}{A} & \etiqueta{RaciI}{I} & \etiqueta{RaciR}{R} & \etiqueta{RaciC}{C} & \textcolor{GrisOscuro}{--} \\
3.6 & Perfil y mis cuartos & \etiqueta{RaciA}{A} & \etiqueta{RaciI}{I} & \etiqueta{RaciR}{R} & \etiqueta{RaciC}{C} & \textcolor{GrisOscuro}{--} \\
\rowcolor{Rama4!15}
\multicolumn{7}{@{}l}{\textbf{\textcolor{Rama4}{4 Pruebas y validación con clientes}}} \\
4.1 & Pruebas del motor & \etiqueta{RaciA}{A} & \etiqueta{RaciI}{I} & \etiqueta{RaciC}{C} & \etiqueta{RaciR}{R} & \textcolor{GrisOscuro}{--} \\
4.2 & Levantamiento de cuartos & \etiqueta{RaciR}{A/R} & \etiqueta{RaciI}{I} & \etiqueta{RaciC}{C} & \etiqueta{RaciI}{I} & \textcolor{GrisOscuro}{--} \\
4.3 & Piloto con clientes & \etiqueta{RaciA}{A} & \etiqueta{RaciR}{R} & \etiqueta{RaciC}{C} & \etiqueta{RaciI}{I} & \textcolor{GrisOscuro}{--} \\
4.4 & Pruebas de aceptación & \etiqueta{RaciA}{A} & \etiqueta{RaciI}{I} & \etiqueta{RaciC}{C} & \etiqueta{RaciR}{R} & \textcolor{GrisOscuro}{--} \\
\rowcolor{Rama5!15}
\multicolumn{7}{@{}l}{\textbf{\textcolor{Rama5}{5 Cierre}}} \\
5.1 & Manuales & \etiqueta{RaciA}{A} & \etiqueta{RaciC}{C} & \etiqueta{RaciR}{R} & \etiqueta{RaciC}{C} & \textcolor{GrisOscuro}{--} \\
5.2 & Informe final y exposición & \etiqueta{RaciR}{A/R} & \etiqueta{RaciC}{C} & \etiqueta{RaciC}{C} & \etiqueta{RaciC}{C} & \etiqueta{RaciI}{I} \\
5.3 & Lecciones aprendidas y acta de cierre & \etiqueta{RaciA}{A} & \etiqueta{RaciC}{C} & \etiqueta{RaciC}{C} & \etiqueta{RaciR}{R} & \etiqueta{RaciC}{C} \\
\end{longtable}
}

\section{Desempeño del Equipo}
El desempeño se gestiona con el ciclo de PM4DEV: Planificar, Hacer, Verificar y Adaptar.

\vspace{0.1cm}
{\small
\rowcolors{2}{white}{GrisFila}
\noindent\begin{tabularx}{\textwidth}{@{}P{2.4cm} L P{2.9cm}@{}}
\toprule
\rowcolor{AzulClaro}
\textbf{Etapa} & \textbf{Qué se hace} & \textbf{Cuándo} \\
\midrule
\textbf{Planificar} & Cada paquete de la EDT se convierte en Issues de GitHub con responsable, fecha y horas estimadas. & Al inicio de cada sprint \\
\textbf{Hacer} & Cada integrante registra sus horas reales y cierra sus Issues con el entregable adjunto. & Durante la semana \\
\textbf{Verificar} & Se comparan las horas y los entregables con lo asignado. Metas: 90\,\% o más de horas cumplidas y 90\,\% o más de entregables a tiempo. & Martes, reunión de control \\
\textbf{Adaptar} & Si un integrante queda por debajo de la meta dos semanas seguidas, se acuerda un apoyo o se reasignan tareas. Al cierre de cada fase, cada integrante evalúa a los demás con una escala de 1 a 5. & Viernes y cierre de fase \\
\bottomrule
\end{tabularx}
}

\section{Incorporación y Liberación}
\begin{itemize}[leftmargin=*]
    \item \textbf{Incorporación:} los cuatro integrantes participan desde la semana 2; no se prevé incorporar personas externas.
    \item \textbf{Reemplazo:} si un integrante se retira, sus paquetes se reasignan según la tabla de habilidades y se actualiza esta matriz.
    \item \textbf{Liberación:} el equipo se libera en la semana 17, después del acta de cierre y del registro de lecciones aprendidas.
\end{itemize}

\section{Firmas}
\noindent En señal de conformidad con el presente plan, los integrantes del equipo firman a continuación:

\vspace{1.5cm}
\begin{center}
\begin{tabular}{c c c}
    \parbox{6.5cm}{\centering \rule{6.0cm}{0.6pt}\\[5pt] \textbf{Clyde Neil Paricahua Pari}\\[2pt] {\footnotesize Project Manager}}
    & \hspace{0.8cm} &
    \parbox{6.5cm}{\centering \rule{6.0cm}{0.6pt}\\[5pt] \textbf{Ruth Karina Apaza Solis}\\[2pt] {\footnotesize Analista de Planificación}}
    \\[2.0cm]
    \parbox{6.5cm}{\centering \rule{6.0cm}{0.6pt}\\[5pt] \textbf{Eliseo Tarqui Ajahuana}\\[2pt] {\footnotesize Especialista Técnico / Arquitectura}}
    & \hspace{0.8cm} &
    \parbox{6.5cm}{\centering \rule{6.0cm}{0.6pt}\\[5pt] \textbf{Viviana Mary Quisbert Quispe}\\[2pt] {\footnotesize Gestora de Comunicación, Riesgos y Documentación}}
\end{tabular}
\end{center}
\end{document}
